% ============================================================================
% Единое Естествознание — Книга III: Инженерный справочник ДКМ
% (book/sources/volumes/vol-engineering-handbook.tex)
% ============================================================================

\chapter*{Введение. Инженерная ветвь единого Естествознания}
\addcontentsline{toc}{chapter}{Введение. Инженерная ветвь единого Естествознания}

Когда физика, химия, биология и физическая реализация вычислений объединяются в единое Естествознание на основе дискретной квантовой механики (ДКМ, $1 + 2^3 = 3^2$, носитель $\mathcal{M} = \mathbb{Z}_{512}$ на оболочке кубооктаэдра $O_2(3)$), электротехника и вычислительная техника сводятся к точному целочисленному согласованию фазовых импедансов. Настоящая \textbf{Книга~III (Инженерный справочник)} опирается на фундаментальную теорию \textbf{Книги~I} (\texttt{formal-foundations-monograph.pdf}) и лабораторные протоколы \textbf{Книги~II} (\texttt{dqm-laboratory-manual.pdf}) и содержит полную систему законов теоретической электротехники фазовой эстафеты (\texttt{LogicCohomology.lean}), физическую реализацию информатики в кремнии без тепла Ландауэра (\texttt{Ring9} / \texttt{Bus36} / \texttt{VerletCell72}) и инженерные расчёты 12-лучевых резонаторов кубооктаэдра $O_2(3)$.

% ============================================================================
\chapter{Теоретическая электротехника фазовой эстафеты на ГЦК-кубооктаэдрах (\texorpdfstring{$V=\mathcal{C}_3,\ I=\mathcal{S}_3$}{V = C3, I = S3})}
\label{ch:eng-electrical-engineering}
% ============================================================================

В классической теории Друде скорость механического дрейфа электронов в медном проводе составляет $v_{\mathrm{drift}} \sim 10^{-4}\text{ м/с}$, что порождает известный парадокс при сопоставлении с передачей электрической мощности вдоль линии со скоростью света. В ДКМ на дискретном носителе $\mathcal{M} = \mathbb{Z}_{512}$ кристаллическая решётка гранецентрированных кубических (ГЦК) проводников ($\mathrm{Cu},\ Z=29$; $\mathrm{Ag},\ Z=47$; $\mathrm{Au},\ Z=79$; $\mathrm{Al},\ Z=13$) образована атомными узлами, чья первая координационная сфера из $12$ ближайших соседей в точности совпадает с кубооктаэдром $O_2(3)$. Электрический ток в таком проводнике представляет собой не механический транспорт шариков, а \textbf{торсионную эстафету вращения фазы} $I(\Delta\varphi) = \mathcal{S}_3(\Delta\varphi)$, сопряжённую с потенциальным натяжением $V(\Delta\varphi) = \mathcal{C}_3(\Delta\varphi) = 128 - |\Delta\varphi|$.

\begin{theorem}[Нагрузочная прямая постоянного тока, согласование импеданса (MPPT) и порог джоулева нагрева, \texttt{LogicCohomology.lean}]
\label{thm:eng-dc-ohm-mppt}
Стационарный сдвиг фазы $\Delta\varphi = k \in [0,128]$ между соседними узлами кубооктаэдра делит бюджет четверти кольца $Q_\square(3) = 128$ между напряжением $V(k) = \mathcal{C}_3(k) = 128 - k$ и током фазовой эстафеты $I(k) = \mathcal{S}_3(k) = k$:
\begin{equation}
  \label{eq:eng-dc-ohm-mppt}
  V(k) + I(k) \;\equiv\; 128, \qquad
  P(k) \;=\; V(k)\,I(k) \;=\; 64^2 - (64 - k)^2 \;\le\; 4\,096 \;=\; K_P(3)^2,
\end{equation}
достигая максимума передачи мощности при $k^* = K_P(3) = 64$ ($V = I = 64$). Внутри подпорогового коридора Гейзенберга $|k| \le 40 < 41$ излучение фононов тождественно равно нулю ($\hat{n}_E(k) = \lfloor|k|/41\rfloor = 0$, баллистический/сверхпроводящий режим); при $|k| \ge 41$ на дефектах решётки излучается $\hat{n}_E(k) \in \{1,\dots,6\}$ фононных квантов (закон Джоуля--Ленца) с рождением вихревых магнитных плакет $B_f = 512\,n_f$.
\end{theorem}

\begin{theorem}[Законы Кирхгофа и индукции Фарадея из нильпотентности кограницы де~Рама, \texttt{LogicCohomology.lean}]
\label{thm:eng-kirchhoff-faraday}
Вокруг всякого замкнутого контура $\partial f$ тождество $(d_1 d_0\varphi)_f \equiv 0$ даёт второй закон Кирхгофа (KVL) в потенциальных цепях и закон электромагнитной индукции Фарадея $(d_1 r)_f = -512(d_1 q)_f = 512\,n_f$ при сцеплении с вихревым магнитным потоком, тогда как $d_2\circ d_1 \equiv 0$ и $\sum_{k=0}^{11}\vec{e}_k = \mathbf{0}$ дают первый закон Кирхгофа (KCL) в каждом узле.
\end{theorem}

\begin{theorem}[Ортогональность реактивной мощности переменного тока и точное гашение в $LC$-резонансе, \texttt{LogicCohomology.lean}]
\label{thm:eng-ac-rlc}
Ёмкостный ($+128$ тиков) и индуктивный ($-128 \equiv +384$ тика) сдвиги фазы на кольце $\mathbb{Z}_{512}$ удовлетворяют
\begin{equation}
  \label{eq:eng-ac-rlc}
  \sum_{\varphi=0}^{511} \mathcal{C}_3(\varphi)\,\mathcal{S}_3(\varphi) = 0, \qquad
  \sum_{\varphi=0}^{511} \mathcal{C}_3(\varphi)^2 = 2\,796\,288, \qquad
  \mathcal{C}_3(\varphi + 128) + \mathcal{C}_3(\varphi + 384) \;\equiv\; 0 \quad (\forall\varphi\in\mathbb{Z}_{512}),
\end{equation}
доказывая равенство нулю работы реактивной мощности за период и поточечную взаимную аннигиляцию напряжений $L$ и $C$ в резонансе.
\end{theorem}

\begin{theorem}[Выпрямители Гретца и бесконденсаторный квадратурный выпрямитель с $0\%$ пульсаций, \texttt{LogicCohomology.lean}]
\label{thm:eng-rectifiers}
Безветвевой однополупериодный выпрямитель $\mathrm{rect}_{1/2}(x) = (|x|+x)/2$ и двухполупериодный 4-диодный мост Гретца $|x|$ дают точные целочисленные средние токи на $\mathbb{Z}_{512}$ без иррациональных множителей $\pi$:
\begin{equation}
  \label{eq:eng-rectifiers}
  \sum_{\varphi=0}^{511}\mathrm{rect}_{1/2}\bigl(\mathcal{S}_3(\varphi)\bigr) = 16\,384 = 128^2 \;\;(\bar{I}_{1/2} = 32), \qquad
  \sum_{\varphi=0}^{511}\bigl|\mathcal{S}_3(\varphi)\bigr| = 32\,768 = 2\times 128^2 \;\;(\bar{I}_{\mathrm{Graetz}} = 64 = K_P),
\end{equation}
а двухфазный квадратурный мостовой выпрямитель ($\varphi$ и $\varphi - 128$) обеспечивает постоянный выход без сглаживающего конденсатора с $0\%$ пульсаций на каждом тике: $|\mathcal{C}_3(\varphi)| + |\mathcal{S}_3(\varphi)| \equiv 128$.
\end{theorem}

\begin{theorem}[Идеальный трансформатор и квадратичное преобразование импеданса]
\label{thm:eng-transformer}
Для первичной и вторичной обмоток $N_1 : N_2$ на общем замкнутом магнитопроводе ($\Phi_{\mathrm{core}} = 512\,n_f$) при $V_1 = N_1\varepsilon_1$, $V_2 = N_2\varepsilon_1$, $I_1 = N_2 i_0$, $I_2 = N_1 i_0$ выполняются $N_2 V_1 - N_1 V_2 \equiv 0$, $N_1 I_1 - N_2 I_2 \equiv 0$, $V_1 I_1 \equiv V_2 I_2$ и $Z_{\mathrm{in}} = (N_1/N_2)^2 Z_{\mathrm{load}}$.
\end{theorem}

% ============================================================================
\chapter{Физическая реализация вычислений без тепла Ландауэра и кремниевая архитектура (\texttt{Ring9} / \texttt{Bus36} / \texttt{VerletCell72})}
\label{ch:eng-silicon-architecture}
% ============================================================================

В то время как абстрактная информатика (полнота Поста и редукция 3-SAT, \texttt{LogicCohomology.lean}) принадлежит чистой математике, её физическая реализация в аппаратуре без внешней ленты Тьюринга определяется тремя инженерными законами ДКМ:
\begin{enumerate}[label=\textbf{\arabic*.},leftmargin=*]
\item \textbf{Подпороговые обратимые вычисления без тепла Ландауэра ($|\Delta\varphi|\le 40 < 41$):}
  двухтактовый симплектический шаг Верле $\varphi(t+1) = (2\varphi(t) - \varphi(t-1) + \Phi_{\mathrm{gate}})\bmod 512$ биективен на 72-битном регистре ячейки (\texttt{VerletCell72}), поэтому любая логическая передача внутри мёртвой зоны Гейзенберга $|\Phi_{\mathrm{gate}}|\le 40$ излучает $\hat{n}_E = 0$ тепловых квантов, обходя тепловой предел стирания Ландауэра $k_B T\ln 2$.
\item \textbf{Диагонализация взаимной индукции между 12 катушками кубооктаэдра $O_2(3)$:}
  вокруг любой выбранной оси катушки $\vec{e}_0\in O_2(3)$ ($\|\vec{e}_0\|_2^2 = 2$) остальные $11$ катушек распадаются на $4$ соседа под углом $60^\circ$ ($\vec{e}_0\cdot\vec{e}_j = +1$), $2$ ортогональных соседа под углом $90^\circ$ ($\vec{e}_0\cdot\vec{e}_j = 0$), $4$ тупоугольных соседа под углом $120^\circ$ ($\vec{e}_0\cdot\vec{e}_j = -1$) и $1$ антиподальный партнёр под углом $180^\circ$ ($\vec{e}_0\cdot(-\vec{e}_0) = -2$), так что внеосевая взаимная индукция сокращается тождественно: $4(+1) + 2(0) + 4(-1) = 0$.
\item \textbf{Безветвевой аппаратный стек (\texttt{hardware/cuboctahedron\_carrier\_asic.v} и \texttt{CarrierCore.cs} на \texttt{.NET~11}):}
  \begin{itemize}[noitemsep]
  \item \texttt{Ring9}: 9-битная модулярная арифметика с бабочкой ротатора Грея $\mathcal{C}_3,\mathcal{S}_3$ и делением Безу;
  \item \texttt{Bus36}: $4\times 9 = 36$-битная четырёхполосная шина по основанию $512$ без ошибок плавающей запятой;
  \item \texttt{VerletCell72}: универсальное 72-битное ядро ячейки (8 девятибитных регистров, 12-соседний затвор и фильтр Паули $D_{\max} = 833\varepsilon$), верифицированное в \url{https://github.com/AI-Guiders/guiders-physics} (16/16 тестов \texttt{xUnit});
  \item \texttt{Module~17 \& 18}: безветвевое восстановление всех $512/512$ фаз из акустического давления $P_{\mathrm{sound}} = \mu_0\mu_1(128-\mathcal{C}_3)$ и хирального тока $J_{\mathrm{rxn}} = \mu_0\mu_1\mathcal{S}_3$, 1-тактовая замкнутая ФАПЧ и 4-канальный нильпотентный томограф Неймана $(2I_4 + N_{\mathrm{nil}})^{-1} = \frac{1}{24}(12 I_4 - 6 N_{\mathrm{nil}} + 3 N_{\mathrm{nil}}^2)$ при $N_{\mathrm{nil}}^3 \equiv 0$.
  \end{itemize}
\end{enumerate}

% applied — том «Прикладная механика кубооктаэдра и фазовые технологии (R&D)» (book/SERIES.md)
\part{Прикладная механика кубооктаэдра, фазовое согласование импеданса и голономные технологии}

\chapter{Масштабная инвариантность кубооктаэдра \texorpdfstring{$O_2$}{O\_2}: что физически разрешено и что запрещено законами сохранения}
\label{ch:applied-scale-invariance-and-bounds}

В томе «Построение» (\bcref{thm:catalan-uniqueness-d3}) было строго установлено и машинно проверено в системе доказательства теорем \textbf{Lean~4} (\texttt{lean/MtCarrier/}, $65$ теорем с нулём аксиом \texttt{sorry}), что уравнение Каталана $1 + 2^d = d^2$ однозначно фиксирует размерность пространства $d = 3$, а три диофантовых замка (делимость Гейзенберга $(1+z)\mid(520+1+z)$, гауссова кулоновская норма $1+8z = 4^2+9^2 = 97$ и замок булеана звезды $2^z = K_P^2 = 4096$) среди всех семи непустых $O_h$-изотропных оболочек в элементарном кубе смещений $\{-1,0,+1\}^3\setminus\{0\}$ однозначно выбирают чистый \textbf{кубооктаэдр гранецентрированной кубической решётки} (\latGCKshort):
\begin{equation}
\label{eq:applied-cuboctahedron-shell}
  O_2 \;=\; \bigl\{e = (e_z, e_y, e_x)\in\{-1,0,+1\}^3 : |e_z| + |e_y| + |e_x| = 2\bigr\},
  \qquad |O_2| = 12,\quad -O_2 = O_2.
\end{equation}

Всякое прикладное и инженерное применение замкнутой теории обязано начинаться с двух вопросов:
\begin{enumerate}[nosep,leftmargin=*]
\item \emph{Почему геометрия кубооктаэдра~$O_2$, выведенная на планковском шаге~$dh$, воспроизводится на атомном, лабораторном и геофизическом масштабах $L = N\,dh$?}
\item \emph{Где проходит точная математическая граница между запрещёнными процессами («вечными двигателями», нарушающими законы сохранения) и разрешёнными параметрическими режимами, недоступными в классических тепловых установках?}
\end{enumerate}

\section{Теорема масштабного самоподобия кубооктаэдра от планковской ячейки до макрорезонатора}
\label{sec:applied-self-similarity}

Ключевое свойство корневой решётки $D_3 = \{(z,y,x)\in\mathbb{Z}^3 : (z+y+x)\bmod 2 = 0\}$ и её первой координационной оболочки~$O_2$ состоит в том, что под действием биномиального оператора макроприбора~$\mathcal{B}_N$ (усреднения по путям длины $2N$ на~$O_2$) тензорная структура операторов переноса и волнового оператора остаётся инвариантной на любом масштабе $L = N\,dh$.

\begin{theorem}[Масштабная инвариантность изотропного лапласиана $\Delta_{O_2(L)}$ и четырёхканального измерения $\mathcal{M}(x,t,N)$]
\label{thm:applied-scale-invariance}
Пусть $O_2(L) := L\cdot O_2 = \{L\,e : e\in O_2\}$ --- масштабированная кубооктаэдральная оболочка с шагом $L = N\,dh$ ($N\ge 1$). Тогда на любом масштабе~$L$ выполняются три точных операторных свойства (\texttt{UnifiedOperators.lean}):
\begin{enumerate}[leftmargin=*]
\item \textbf{Инвариантность тензора второго момента и ППН-тензора (\texttt{cuboctahedron\_\allowbreak second\_\allowbreak moment\_\allowbreak isotropy\_\allowbreak tensor}):}
  для всех $e\in O_2(L)$ квадрат длины ребра равен $\|e\|_2^2 = 2L^2$, матрица вторых моментов строго диагональна и изотропна:
  \begin{equation}
  \label{eq:applied-second-moment-L}
    \sum_{e\in O_2(L)} e_i\,e_j \;=\; 8L^2\,\delta_{ij} \;=\; |\triangle(3)|\,L^2\,\delta_{ij},
    \qquad
    \sum_{e\in O_2(L)} \frac{e_i\,e_j}{\|e\|_2^2} \;=\; 4\,\delta_{ij} \;=\; d_{\mathrm{in}}(\square)\,\delta_{ij},
  \end{equation}
  а дискретный лапласиан $(\Delta_{O_2(L)} f)(x) = \sum_{e\in O_2(L)}(f(x+e)-f(x))$ на любом трёхмерном квадратичном поле $f(z,y,x) = ax^2+by^2+cz^2+dxy+exz+gyz+px+qy+rz+u$ тождественно аннулирует все перекрёстные члены $xy,xz,yz$ и возвращает $8L^2(a+b+c)$ (\texttt{laplacian\_\allowbreak O2\_\allowbreak exact\_\allowbreak on\_\allowbreak quadratic\_\allowbreak polynomial}).
\item \textbf{Точное сохранение калибровочных и топологических законов на любом окне прибора $N\ge 1$ (\texttt{HolonomyMaxwell.lean}):}
  под действием биномиального окна прибора $\mathcal{B}_N$ квантование плакетного потока $(d_1\mathcal{B}_N[A])_f = 512\,\mathcal{B}_N[n_f]$ (\texttt{instrument\_\allowbreak convolution\_\allowbreak plaquette\_\allowbreak flux\_\allowbreak exact}), закон отсутствия изолированных монополей $\operatorname{div}_{\mathrm{cube}}\mathcal{B}_N[n_f] \equiv 0$ (\texttt{instrument\_\allowbreak convolution\_\allowbreak gauss\_\allowbreak monopole\_\allowbreak free\_\allowbreak exact}) и закон индукции Фарадея $\Delta_t B_f^{(N)} + (d_1 E^{(N)})_f \equiv 0$ (\texttt{instrument\_\allowbreak convolution\_\allowbreak faraday\_\allowbreak law\_\allowbreak exact}) выполняются с нулевой погрешностью при любом конечном~$N\ge 1$.
\item \textbf{Иерархия физических реализаций кубооктаэдра $O_2(L)$:}
  поскольку свободная энергия Гинзбурга--Ландау--Бома достигает минимума на изотропной оболочке~\eqref{eq:applied-second-moment-L}, одна и та же $12$-вершинная геометрия~$O_2(L)$ реализуется на четырёх макроскопических этажах:
  \begin{itemize}[nosep,leftmargin=*]
  \item \emph{Кристаллический этаж ($L \sim 2{,}5\text{--}4\ \text{\AA}$, $N\sim 10^{25}$):} атомные ГЦК-решётки переходных и благородных металлов ($\mathrm{Ni},\ \mathrm{Pd},\ \mathrm{Pt},\ \mathrm{Cu},\ \mathrm{Ag},\ \mathrm{Au},\ \mathrm{Al},\ \gamma\text{-}\mathrm{Fe}$), в которых каждый узел окружён $12$ ближайшими соседями по вершинам кубооктаэдра;
  \item \emph{Мезоскопический и кавитационный этаж ($L \sim 1\ \text{мкм}\text{--}1\ \text{мм}$, $N\sim 10^{29}\text{--}10^{32}$):} стоячие акустические и плазмонные моды в сферических и додекаэдральных микрополостях;
  \item \emph{Лабораторный резонаторный этаж ($L \sim 1\ \text{см}\text{--}1\ \text{м}$, $N\sim 10^{33}\text{--}10^{35}$):} $12$-лучевые фазированные акустико-электромагнитные концентраторы с расположением излучателей в вершинах~$O_2(L)$;
  \item \emph{Планетарный геофизический этаж ($L \sim 10^3\text{--}10^4\ \text{км}$, $N\sim 10^{41}\text{--}10^{42}$):} внутреннее $\mathrm{Fe}\text{-}\mathrm{Ni}$ ядро и мантия планет земной группы.
  \end{itemize}
\end{enumerate}
\end{theorem}

\section{Жёсткая демаркация: что запрещено законами сохранения, а что разрешено}
\label{sec:applied-permitted-vs-forbidden}

Всякая инженерная программа, опирающаяся на дискретную модель носителя, должна прежде всего провести строгую границу между невозможным и возможным:

\begin{table}[htbp]
\centering
\footnotesize
\setlength{\tabcolsep}{4pt}
\caption{Демаркация запрещённых и разрешённых процессов в целочисленной механике кубооктаэдра}
\label{tab:applied-demarcation}
\begin{tabular}{@{}p{0.28\linewidth}p{0.33\linewidth}p{0.33\linewidth}@{}}
\toprule
Физический вопрос & Что строго запрещено (и почему) & Что строго разрешено (и какой теоремой доказано) \\
\midrule
Рождение энергии «из ничего» / нарушение нормы & Изменение полной квадратичной нормы ячейки $\mu_0^2+\mu_1^2$ при свободном шаге (\texttt{polar\_\allowbreak step\_\allowbreak uniqueness\_\allowbreak from\_\allowbreak a3\_\allowbreak a4\_\allowbreak a5}: $\mu_c'=\mu_c$) & Перевод запаса фазовой связности $B_f = 512\,n_f$ в излучение $A_e$ при парной перестройке или аннигиляции вихрей $(+1,-1)\to(0,0)$ (\texttt{shared\_\allowbreak edge\_\allowbreak charge\_\allowbreak conservation}) \\
\addlinespace
Рождение одиночного топологического/магнитного заряда & Появление изолированного монополя $\sum_{f\in\partial c} n_f \ne 0$ (\texttt{discrete\_\allowbreak gauss\_\allowbreak cube\_\allowbreak divergence\_\allowbreak zero}: $\operatorname{div}_{\mathrm{cube}} n_f \equiv 0$) & Рождение замкнутых нейтральных петель и диполей $(+1,-1)$ или трёхгранных узлов $K_3$ $(+1,+1,-2)$ при переходе ребра через $|A_e|=256$ \\
\addlinespace
Преодоление кулоновского/упругого барьера слияния вихрей & Беспороговое слияние при закрытом затворе Гейзенберга $|\Phi|<41$ (\texttt{a6\_\allowbreak a8\_\allowbreak gate\_\allowbreak pauli\_\allowbreak ceiling\_\allowbreak 256\_\allowbreak and\_\allowbreak heisenberg\_\allowbreak deadband\_\allowbreak 41}) & Когерентное снижение знаменателя барьера $D(x) = \zeta_{\mathrm{real}}+128\mu^2+8192$ в $25\text{--}587$ раз до пола $\varepsilon=8192$ при противофазной накачке $12$ лучей~$O_2$ (\texttt{coherent\_\allowbreak cuboctahedron\_\allowbreak barrier\_\allowbreak collapse\_\allowbreak 587x}) \\
\addlinespace
Предсказание локальной геологии без краевых замеров & Вычисление текущих координат залежей из одного $t=0$ без датчиков (неизвестен $60$-значный планковский адрес и пройдено дерево ветвей Паули) & Замкнутая $4$-канальная голономная томография профиля импеданса $D(y)$ внутри объёма по краевому отклику $\mathcal{M}(x,t,N)$ (\texttt{closed\_\allowbreak form\_\allowbreak unified\_\allowbreak 4channel\_\allowbreak measurement\_\allowbreak iv263\_\allowbreak t2}) \\
\bottomrule
\end{tabular}
\end{table}

\chapter{Когерентное схлопывание импеданса барьера в 12-лучевом кубооктаэдре: решёточный синтез против теплового нагрева}
\label{ch:applied-coherent-barrier-collapse}

Рассмотрим фундаментальное уравнение фазовых ворот~\latGCKshort\ (\texttt{evalGateFromHolonomy} в \texttt{TrajectoryIV263.lean}):
\begin{equation}
\label{eq:applied-gate-law}
  \Phi(x,t) \;=\; \mathrm{wrap}_{512}\!\left(\left\lfloor \frac{64\,\bigl(\zeta_{\mathrm{imag}}(x) + \operatorname{sgn}(\zeta_{\mathrm{imag}}(x))\,\delta\zeta_{\mathrm{real}}(x)\bigr)}{D(x)} \right\rfloor\right),
\end{equation}
где числитель управляется нечётным током переноса $\zeta_{\mathrm{imag}}(x) = \mu(x)\sum_{y\in N_{12}(x)}\mu(y)\,S(A_{xy})$ и чётным дефицитом Бома $\delta\zeta_{\mathrm{real}}(x) = \mu(x)\sum_{y\in N_{12}(x)}\mu(y)\,|A_{xy}| \ge 0$ (\texttt{even\_\allowbreak holonomy\_\allowbreak bohm\_\allowbreak deficit\_\allowbreak eq\_\allowbreak abs\_\allowbreak connection}), а знаменатель~$D(x)$ представляет собой реактивный импеданс (барьер жёсткости) узла:
\begin{equation}
\label{eq:applied-denominator-impedance}
  D_{\mathrm{raw}}(x) \;=\; \zeta_{\mathrm{real}}(x) + 128\,\mu(x)^2 + \varepsilon,
  \qquad
  D(x) \;=\; \begin{cases}
    D_{\mathrm{raw}}(x), & \text{если } D_{\mathrm{raw}}(x) > 0,\\
    \varepsilon = 8\,192, & \text{если } D_{\mathrm{raw}}(x) \le 0,
  \end{cases}
\end{equation}
где $\zeta_{\mathrm{real}}(x) = \mu(x)\sum_{y\in N_{12}(x)}\mu(y)\,C(A_{xy}) = \mu(x)\sum_{y\in N_{12}(x)}\mu(y)\,\bigl(128 - |A_{xy}|\bigr)$.

Для того чтобы в узле~$x$ произошла топологическая перестройка (слияние двух вихревых ядер или рождение новой связанной пары через разрез ветви Паули $|A_{xy}| = 256$), фазовый удар ворот~$\Phi(x,t)$ должен преодолеть порог Гейзенберга $|\Phi| \ge 41$ и вывести разность фаз на ребре за предел полуоборота Паули $\Phi_{\mathrm{Pauli}} = 256$ тиков.

Из структуры дроби~\eqref{eq:applied-gate-law} непосредственно видно, что существует \textbf{два физически противоположных способа} достичь $|\Phi| \sim 256$:

\section{Тепловой (некогерентный) режим против когерентного противофазного режима второго квадранта}
\label{sec:applied-thermal-vs-coherent}

\begin{theorem}[Когерентное схлопывание барьера кубооктаэдра в $25\text{--}587$ раз]
\label{thm:applied-barrier-collapse}
Пусть центральный узел~$x$ и его $12$ соседей по кубооктаэдру~$O_2$ имеют амплитуду $\mu \in \{1,\dots,64\}$ ($\varepsilon = 64\times 128 = 8\,192$, \texttt{coherent\_\allowbreak cuboctahedron\_\allowbreak barrier\_\allowbreak collapse\_\allowbreak 587x}). Тогда:
\begin{enumerate}[leftmargin=*]
\item \textbf{Некогерентный / тепловой режим первого квадранта ($|A_{xy}| \le 41$, $C(A_{xy}) \ge 87$):}
  когда относительные фазы соседей лежат в первом квадранте на границе порога Гейзенберга ($|A_{xy}| = 41 \implies C(41) = 128 - 41 = +87 > 0$), продольная голономия положительна и складывается с собственным барьером узла $+128\,\mu^2$:
  \begin{equation}
  \label{eq:applied-d-thermal}
    D_{\mathrm{thermal}}(\mu) \;=\; 12\cdot 87\,\mu^2 + 128\,\mu^2 + 8\,192 \;=\; 1\,172\,\mu^2 + 8\,192.
  \end{equation}
  При амплитуде нуклеации $\mu = 13$ импеданс барьера равен $D_{\mathrm{thermal}}(13) = 206\,260$, а на потолке плотности $\mu = 64$ достигает $D_{\mathrm{thermal}}(64) = 4\,808\,704$ (в полном синфазном пределе $A_{xy}=0,\ C(0)=128$: $D_{\max}(64) = 13\cdot 128\cdot 64^2 + 8\,192 = 6\,823\,936$). Преодоление такого знаменателя требует гигантской накачки числителя~$z_{\mathrm{drive}}$ (тепловой пробой / горячий термояд при $T \sim 10^8\ \text{К}$).
\item \textbf{Когерентный противофазный режим второго квадранта ($139 \le |A_{xy}| \le 256$, $C(A_{xy}) \le -11$):}
  если все $12$ направлений кубооктаэдра~$O_2$ вокруг узла~$x$ возбуждены когерентной стоячей волной во втором квадранте фазового кольца $|A_{xy}| \ge 139$ тиков ($\Delta\varphi \ge 97{,}7^\circ$), то каждая из $12$ граневых проекций становится отрицательной:
  \[
    C(A_{xy}) \;=\; 128 - |A_{xy}| \;\le\; 128 - 139 \;=\; -11 \;<\; 0.
  \]
  Тогда продольная голономия $\zeta_{\mathrm{real}}(x) \le 12\cdot(-11)\mu^2 = -132\,\mu^2$ \textbf{полностью компенсирует собственный барьер узла $+128\,\mu^2$}:
  \begin{equation}
  \label{eq:applied-d-coherent}
    D_{\mathrm{raw,coh}}(\mu) \;\le\; -132\,\mu^2 + 128\,\mu^2 + 8\,192 \;=\; 8\,192 - 4\,\mu^2,
  \end{equation}
  и при $\mu = 64$ ($D_{\mathrm{raw,coh}}(64) = -8\,192 \le 0$) знаменатель ворот~\eqref{eq:applied-denominator-impedance} схлопывается до минимального вакуумного пола:
  \[
    D_{\mathrm{coherent}}(64) \;=\; \varepsilon \;=\; 8\,192.
  \]
\item \textbf{Коэффициент параметрического просветления барьера:}
  отношение теплового импеданса барьера к когерентному импедансу в фокусе кубооктаэдра составляет:
  \begin{equation}
  \label{eq:applied-barrier-ratio}
  \begin{aligned}
    \left\lfloor\frac{D_{\mathrm{thermal}}(13)}{\varepsilon}\right\rfloor &= \left\lfloor\frac{206\,260}{8\,192}\right\rfloor = \mathbf{25}, \\
    \frac{D_{\mathrm{thermal}}(64)}{\varepsilon} &= \frac{4\,808\,704}{8\,192} = \mathbf{587},
    \qquad
    \frac{D_{\max}(64)}{\varepsilon} = \frac{6\,823\,936}{8\,192} = \mathbf{833}.
  \end{aligned}
  \end{equation}
\end{enumerate}
\end{theorem}

\begin{proof}
Все равенства~\eqref{eq:applied-d-thermal}--\eqref{eq:applied-barrier-ratio} следуют из прямой целочисленной подстановки $C(41) = 128 - 41 = 87$, $C(0) = 128$ и $C(139) = 128 - 139 = -11$ в~\eqref{eq:applied-denominator-impedance}; машинное доказательство без аксиом содержится в теореме \texttt{coherent\_\allowbreak cuboctahedron\_\allowbreak barrier\_\allowbreak collapse\_\allowbreak 587x} (\texttt{UnifiedOperators.lean}), а необходимость самого защитного пола $D\le 0\mapsto\varepsilon=8\,192$ для рождения вихревой пары на такте $t=1\to 2$ доказана теоремой-абляцией \texttt{ablation\_\allowbreak no\_\allowbreak denom\_\allowbreak guard\_\allowbreak corrupts\_\allowbreak t2\_\allowbreak vortex\_\allowbreak pair} (\texttt{TrajectoryIV263.lean}).
\end{proof}

\section{Технический протокол 12-лучевого кубооктаэдрального фазового концентратора}
\label{sec:applied-12beam-protocol}

Теорема~\ref{thm:applied-barrier-collapse} объясняет, почему низкоэнергетические ядерные и вихревые реакции в ГЦК-кристаллах ($\mathrm{Pd},\ \mathrm{Ni},\ \mathrm{Cu}$), насыщенных дейтерием или водородом, случайно и нестабильно наблюдались в отдельных микротрещинах при электролизе или ударном нагреве, но не воспроизводились при стационарном однородном насыщении:
в равновесной решётке $|A_{xy}| < 128$, проекции $C(A_{xy}) > 0$ и знаменатель барьера~$D$ остаётся высоким.

Для детерминированного воспроизведения режима $D(x) \to \varepsilon = 8\,192$ установка должна иметь следующую архитектуру:
\begin{enumerate}[leftmargin=*]
\item \textbf{Геометрия 12-лучевой камеры ($O_2(L)$):}
  рабочая камера выполняется в форме сферы или кубооктаэдра, в $12$ вершинах которого $(\pm L, \pm L, 0),\ (\pm L, 0, \pm L),\ (0, \pm L, \pm L)$ ($6$ коллинеарных антиподальных пар $\pm e$) установлены $12$ идентичных фазовых излучателей (ультразвуковых/кавитационных пьезодрайверов и синхронных ВЧ/СВЧ-индукторов).
\item \textbf{Ориентация мишени (\latGCKshort-кристалл):}
  в геометрическом центре фокуса размещается монокристаллическая или текстурированная мишень из ГЦК-металла ($\mathrm{Ni}$ или $\mathrm{Pd}$), чьи кристаллографические оси $[100],[010],[001]$ совмещены с осями внешнего $12$-лучевого кубооктаэдра.
\item \textbf{Девятибитный ($512$-шаговый) синтез фазы и маска $\mathrm{IV}(3) = 263$:}
  многоканальный цифровой синтезатор (DDS/FPGA) квантует период несущей частоты на $\phaseRingCount = 512$ фазовых тиков и формирует сходящийся двухтактовый пакет:
  \begin{itemize}[nosep,leftmargin=*]
  \item на первом такте ($t=0\to 1$) на $12$ излучателей подаётся фазовая маска трёхгранного барионного угла~$K_3$ ($\mathrm{IV}(3) = 263 = \texttt{0b1\_00000111}$: сдвиг активных октантов на $3\times 41 = 123$ тика $\approx 86{,}48^\circ$), создающая на трёх координатных доменных стенках встречный фазовый удар $+76$ и $-77$ тиков (\texttt{iv263\_flat\_domain\_wall\_asymmetry\_t0});
  \item к моменту фокусировки ($t=1$) суперпозиция прямой волны и инерции второго порядка $2\varphi(1)-\varphi(0)$ переводит разность фаз в центральной зоне во второй квадрант $|A_{xy}| \in [139, 256]$ тиков ($C(A_{xy}) < 0$), обнуляя эффективный знаменатель барьера до $D = \varepsilon = 8\,192$ и вызывая на такте $t=2$ когерентный переброс через разрез Паули $|A_e| = 256$.
  \end{itemize}
\end{enumerate}

\chapter{Квантовый фазовый храповик на доменной стенке и топологический аккумулятор \texorpdfstring{$B_f = 512\,n_f$}{B\_f = 512 n\_f}}
\label{ch:applied-phase-ratchet-and-vortex-battery}

Помимо параметрического снижения барьера~$D(x)$, точная целочисленная механика ворот~\eqref{eq:applied-gate-law} и деления с остатком $(d_0\varphi)_e = 512\,k_e + A_e$ открывает два твердотельных энергетических механизма, которые тождественно исчезают в гладком континуальном приближении.

\section{Киральный выпрямитель (храповик $-1$ тик) на целочисленном делении вниз \texttt{Int.fdiv}}
\label{sec:applied-chiral-ratchet}

В теореме \texttt{iv263\_flat\_domain\_wall\_asymmetry\_t0} и теореме-абляции \texttt{ablation\_\allowbreak tdiv\_\allowbreak destroys\_\allowbreak domain\_\allowbreak wall\_\allowbreak chiral\_\allowbreak asymmetry} (\texttt{TrajectoryIV263.lean}) доказано фундаментальное свойство целочисленного деления с остатком вниз $\lfloor u / D\rfloor$:
для любого положительного числителя $u > 0$, не кратного знаменателю $D > 0$ ($u \bmod D \ne 0$), выполняется точное тождество
\begin{equation}
\label{eq:applied-floor-ratchet}
  \left\lfloor \frac{+u}{D} \right\rfloor + \left\lfloor \frac{-u}{D} \right\rfloor \;=\; -1.
\end{equation}
На плоской доменной стенке кристалла $\mathrm{IV}(3)=263$ ($u = 64\times 195\,702 = 12\,524\,928$, $D = 164\,686$) положительный берег стенки получает импульс $\Phi_+ = \lfloor +12\,524\,928 / 164\,686\rfloor = +76$ тиков, а отрицательный берег --- импульс $\Phi_- = \lfloor -12\,524\,928 / 164\,686\rfloor = -77$ тиков, так что суммарный сдвиг через стенку равен $\Phi_+ + \Phi_- = -1$ тику за такт.

\textbf{Техническая реализация (твердотельный фазовый выпрямитель):}
если в гетероструктуре или многослойном ГЦК-переходе с закреплённой доменной стенкой ($\Delta\varphi_0 = 123$ тика) возбуждать чисто симметричные знакопеременные колебания тока $\pm\zeta_{\mathrm{imag}}$ с амплитудой выше порога Гейзенберга ($|\Phi| \ge 41$, так что $u \bmod D \ne 0$), то на каждом периоде качания асимметрия~\eqref{eq:applied-floor-ratchet} накапливает однонаправленный фазовый сдвиг $\Delta\varphi_{\mathrm{net}} = -1$ тик за цикл без внешнего постоянного смещения, преобразуя симметричный высокочастотный шум или акустическую накачку в направленную фазовую ЭДС $(d_0\Phi)_e$.

\section{Топологический аккумулятор и квантованный сброс \texorpdfstring{$512$}{512} тиков при рекомбинации вихревых петель}
\label{sec:applied-topological-battery}

По теореме~\bcref{thm:discrete-maxwell-sourceless} (\texttt{plaquette\_\allowbreak flux\_\allowbreak eq\_\allowbreak 512\_\allowbreak mul\_\allowbreak winding} в \texttt{HolonomyMaxwell.lean}) каждый элементарный вихревой плакет $n_f = \pm 1$ удерживает в своей границе строго квантованную циркуляцию реберной связности:
\begin{equation}
\label{eq:applied-plaquette-512-storage}
  B_f \;=\; (d_1 A)_f \;=\; A_{01} + A_{12} + A_{23} + A_{30} \;=\; 512\,n_f.
\end{equation}
Пока все четыре ребра плакета находятся вдали от разреза Паули $|A_e| = 256$, целочисленный заряд $n_f \in \mathbb{Z}$ топологически защищён: никакие малые возмущения $|A_e| < 256$ не могут изменить $n_f$ (\texttt{plaquette\_\allowbreak winding\_\allowbreak zero\_\allowbreak of\_\allowbreak principal}).
При управляемом сведении двух противоположно ориентированных плакетов $(n_{f_+}=+1,\ n_{f_-}=-1)$ на общее ребро~$e_{\mathrm{shared}}$ и подаче импульса, переводящего $A_{e_{\mathrm{shared}}}$ обратно через порог $\pm 256$ ($\Delta k_{e_{\mathrm{shared}}} = \pm 1$), пара аннигилирует ($(+1,-1)\to(0,0)$ по теореме \texttt{shared\_edge\_charge\_conservation}), высвобождая запасённую фазовую циркуляцию $512$ тиков напрямую в направленный импульс электрического поля $E_e = -\Delta_t A_e - (d_0\Phi)_e$ (\texttt{instrument\_\allowbreak convolution\_\allowbreak faraday\_\allowbreak law\_\allowbreak exact}).

\chapter{Четырёхканальная голономная томография сред и фазовое геодинамо ядер планет}
\label{ch:applied-tomography-and-planetary-cores}

Замкнутая формула измерения макроприбора~$\mathcal{B}_N$ (\bcref{thm:continuum-emergence-from-z}(iii), \texttt{closed\_\allowbreak form\_\allowbreak unified\_\allowbreak 4channel\_\allowbreak measurement\_\allowbreak iv263\_\allowbreak t2}) связывает все четыре физических канала в одной точке пространства-времени $(x,t)$ через одно и то же ядро путей кубооктаэдра $W_N(x-y)$ и один и тот же знаменатель импеданса $D(y,t) = \zeta_{\mathrm{real}}(y,t) + 128\mu(y)^2 + 8\,192$.

\section{Безынтегральная голономная дефектоскопия и геофизическая томография}
\label{sec:applied-holonomic-tomography}

В классической геофизике и неразрушающем контроле акустическая сейсмика, электромагнитное зондирование и гравиметрия решаются как три несвязанные обратные задачи с эмпирическими подгоночными коэффициентами породы.
В замкнутой механике кубооктаэдра все четыре отклика зондирующего пакета на неоднородность $\mu(y)$ определяются единой алгебраической матрицей:
\begin{enumerate}[leftmargin=*]
\item \textbf{Нечётный акустический сдвиг тока:} $J_{xy} = \mu_x\mu_y S(A_{xy})$ (попарно консервативный на всех $512$ углах, \texttt{odd\_\allowbreak holonomy\_\allowbreak current\_\allowbreak pairwise\_\allowbreak conservation\_\allowbreak all\_\allowbreak 512});
\item \textbf{Чётный фазовый дефицит Бома (дисперсионный скос пакета):}
  \begin{equation}
  \label{eq:applied-bohm-tomography}
    \delta K_{xy} \;=\; 128\,\mu_x\mu_y - \mu_x\mu_y C(A_{xy}) \;=\; \mu_x\mu_y\,|A_{xy}| \;\ge\; 0
  \end{equation}
  (\texttt{even\_\allowbreak holonomy\_\allowbreak bohm\_\allowbreak deficit\_\allowbreak eq\_\allowbreak abs\_\allowbreak connection});
\item \textbf{Задержка времени пробега (оптический показатель преломления):}
  \begin{equation}
  \label{eq:applied-refraction-tomography}
    n_{\mathrm{opt}}^{(N)}(x,t) \;=\; \frac{1}{D_{\mathrm{vac}}}\sum_{y} W_N(x-y)\,\bigl(\zeta_{\mathrm{real}}(y,t) + 128\mu(y)^2 + 8\,192\bigr) \;>\; 0
  \end{equation}
  (\texttt{gravity\_\allowbreak denominator\_\allowbreak strictly\_\allowbreak positive\_\allowbreak floor\_\allowbreak 8192}).
\end{enumerate}
Поскольку для каждого ребра выполняется точное линейное соотношение $C(A_{xy}) + |A_{xy}| \equiv 128$, одновременное измерение на поверхности приёмной антенной решёткой (размещённой по узлам плоского гексагонального среза $\{111\}$ кубооктаэдра) времени задержки~\eqref{eq:applied-refraction-tomography} и дисперсионного дефицита Бома~\eqref{eq:applied-bohm-tomography} позволяет исключить неизвестную фазу $|A_{xy}|$ и напрямую восстановить карту квадрата плотности $\mu(y)^2$ внутри объёма без решения некорректных интегральных уравнений Фредгольма.

\section{Строение ядер планет и механизм фазового геодинамо на пороге Паули \texorpdfstring{$A_{16}$}{A16}}
\label{sec:applied-planetary-cores}

Применим замкнутые операторы плотностного потолка~\Axref{ax:A7} ($\rho_q = \mu_0^2+\mu_1^2 \le 4096$), оптического преломления~\eqref{eq:applied-refraction-tomography} и спинорного исключения Паули~\Axref{ax:A16} (\texttt{pauliPredicate} в \texttt{TrajectoryIV263.lean}) к гравитационному сжатию вещества в центральном ядре планеты:
\begin{enumerate}[leftmargin=*]
\item \textbf{Формирование ГЦК-ядра и гравитационная фокусировка:}
  под действием реберного преломления $\Delta_t A_\perp = (d_0\Phi)_\perp < 0$ в сторону роста знаменателя~$D(x)$ тяжёлые элементы ($\mathrm{Fe},\ \mathrm{Ni}$) концентрируются к центру масс, где при высоком давлении кристаллизуются в плотнейшую гранецентрированную кубическую фазу ($\gamma\text{-}\mathrm{Fe}/\mathrm{Ni}$ с координационным кубооктаэдром~$O_2$, $|N_{12}|=12$).
\item \textbf{Включение затвора Паули $A_{16}$ на половине потолка плотности ($\mu_c^2 > 2048$) и генерация магнитного поля:}
  когда в сжатом внутреннем ядре квадраты спинорных амплитуд превышают половину планковского потолка ($\mu_0^2 > K_P^2/2 = 2048$ и $\mu_1^2 > 2048$, то есть $\mu_c \ge 46$ квантов из $64$) при сонаправленных фазах $|\varphi_1-\varphi_0| < 41$, срабатывает предикат Паули~\Axref{ax:A16}, вносящий полуоборот $\Phi_{\mathrm{Pauli}} = 256$ тиков ($\pi\ \text{рад}$).
\item \textbf{Циклический сброс магнитного потока $B_f = 512\,n_f$ и инверсии геодинамо:}
  скачок фазы на $256$ тиков на границе внутреннего и внешнего ядра перебрасывает реберные разности $(d_0\varphi)_e$ через разрез $\pm 256$, порождая по теореме~\bcref{thm:discrete-maxwell-sourceless} макроскопические вихревые кольца магнитного потока $B_f^{(N)} = 512\,\mathcal{B}_N[n_f]$ вдоль осей кубооктаэдра (при строго нулевой дивергенции $\operatorname{div}_{\mathrm{cube}} B^{(N)} \equiv 0$, то есть чисто дипольно-квадрупольной структуре без магнитных монополей) и одновременно меняя знак продольной голономии $C(256) = -128 < 0$, что ограничивает дальнейшее сжатие ядра и порождает квазипериодические циклы геомагнитных экскурсов и инверсий.
\end{enumerate}

\input{volumes/include-computational-laboratory.tex}

\input{volumes/include-measurement-processing.tex}

\chapter{Критические эксперименты и инженерные гейты верификации (серия CE-RD)}
\label{ch:applied-critical-experiments}

\begin{labseries}{Серия CE-RD: когерентное схлопывание барьера кубооктаэдра, фазовый храповик и 4-канальная голономная томография}{sec:lab-series-ce-rd}
\begin{labrecord}{CE-RD-01 · Когерентное схлопывание импеданса барьера \texorpdfstring{$D(x)\to\varepsilon=8\,192$}{D(x)->8192} в 12-лучевом кубооктаэдре (\texorpdfstring{$25\times\text{--}587\times$}{25x-587x})}{lab:ce-rd-01}
\label{exp:cerd01}

\subsubsection*{Постановка}
\textbf{Физический вопрос.}
Во сколько раз когерентная противофазная накачка $12$ вершин кубооктаэдра~$O_2$ во втором квадранте фазового кольца ($139 \le |A_{xy}| \le 256$, $C(A_{xy}) \le -11$) снижает реактивный знаменатель барьера $D(x) = \zeta_{\mathrm{real}}(x) + 128\mu(x)^2 + \varepsilon$ по сравнению с некогерентным тепловым режимом ($|A_{xy}| \le 41$, $C(A_{xy}) \ge +87$), и разрушается ли рождение вихревой пары на такте $t=2$ при отключении защитного пола $D \le 0 \mapsto \varepsilon = 8\,192$?

\textbf{Что проверяем.}
\begin{enumerate}[label=\textbf{(\alph*)},leftmargin=*,itemsep=0.25em]
\item Точные целочисленные значения импеданса $D_{\mathrm{thermal}}(13) = 206\,260$, $D_{\mathrm{thermal}}(64) = 4\,808\,704$ и $D_{\max}(64) = 6\,823\,936$ в первом квадранте (\texttt{coherent\_\allowbreak cuboctahedron\_\allowbreak barrier\_\allowbreak collapse\_\allowbreak 587x}).
\item Схлопывание знаменателя $D_{\mathrm{coherent}}(64) = \varepsilon = 8\,192$ при $|A_{xy}| = 139$ ($C(139) = -11$, $D_{\mathrm{raw}} = -8\,192 \le 0$) и отношения барьеров $25\times$, $587\times$ и $833\times$.
\item Реальное срабатывание порога $D_{\mathrm{raw}} = -113\,826 \le 0 \mapsto D = 8\,192$ в узле $(-2,-1,-1)$ на такте $t=1$ эволюции $\mathrm{IV}(3)=263$ (\texttt{ablation\_\allowbreak no\_\allowbreak denom\_\allowbreak guard\_\allowbreak corrupts\_\allowbreak t2\_\allowbreak vortex\_\allowbreak pair}).
\end{enumerate}

\textbf{Рабочая гипотеза.}
Когерентная противофазная накачка по $12$ направлениям кубооктаэдра снижает знаменатель барьера ровно в $25$ раз при $\mu=13$ и в $587$ раз при $\mu=64$ ($833$ раза относительно синфазного потолка), а снятие защитного пола $D>0$ роняет число рождённых $xy$-вихрей на $t=2$ с $26$ до $24$.

\textbf{Наблюдаемые.}
Знаменатели $D_{\mathrm{thermal}}(\mu)$, $D_{\mathrm{coherent}}(\mu)$, коэффициенты редукции барьера $K_{13}, K_{64}, K_{\max}$ и число рождённых $xy$-вихрей $N_{xy}(t=2)$ (табл.~\ref{tab:exp-ce-rd-01}).

\textbf{Критерий согласия с теорией.}
Точные целочисленные равенства $K_{13} = 25$, $K_{64} = 587$, $K_{\max} = 833$ и $N_{xy}(t=2) = 26$ (тип~B, машинная верификация в Lean~4).

\subsubsection*{Теоретическая часть}
По теореме~\ref{thm:applied-barrier-collapse} при $|A_{xy}| = 139$ тиков проекция грани равна $C(139) = 128 - 139 = -11$, так что сумма по $12$ соседям кубооктаэдра $\zeta_{\mathrm{real}} = -132\mu^2$ компенсирует собственный барьер $+128\mu^2$, переводя знаменатель~$D$ на вакуумный пол $\varepsilon = 8\,192$.

\textbf{Проверка размерностей.}
В естественных единицах ядра~$\Mlayer$ амплитуда $\mu\in\{1,\dots,64\}$ измеряется в квантах, проекция грани $C(A)\in[-128,128]$ --- в долях $Q_\square=128$, а все слагаемые знаменателя~$D$ ($\zeta_{\mathrm{real}}$, $128\mu^2$ и $\varepsilon = 64\times 128 = 8\,192$) имеют единую размерность $\mathrm{кванты}^2\cdot Q_\square$, поэтому отношение $K = D_{\mathrm{thermal}}/D_{\mathrm{coherent}}$ строго безразмерно ($[K] = 1$).

\subsubsection*{Ход эксперимента}
{\sloppy Воспроизведение: \texttt{lake build} в \texttt{lean/} (\texttt{UnifiedOperators.lean}, \texttt{TrajectoryIV263.lean}) и \texttt{python verify\_principles.py --suite m\_axioms}.\par}

\begin{enumerate}[label=\textbf{Шаг \arabic*:},leftmargin=*,itemsep=0.35em]
\item Вычислить в целочисленной арифметике $D_{\mathrm{thermal}}(\mu)$ при $|A_{xy}|=41$ и $D_{\mathrm{coherent}}(\mu)$ при $|A_{xy}|=139$ для $\mu\in\{13, 64\}$.
\item Проверить срабатывание условия $D_{\mathrm{raw}} \le 0 \mapsto \varepsilon = 8\,192$ на такте $t=1$ траектории $g^t(\mathrm{IV}(3)=263)$ в узле $(-2,-1,-1)$.
\item Сравнить полный 3D-ценз вихревых плакетов при $t=2$ с защитным полом $D>0$ и при его отключении.
\end{enumerate}

\begin{table}[htbp]
\centering
\footnotesize
\setlength{\tabcolsep}{4pt}
\caption{Верификация когерентного схлопывания барьера кубооктаэдра (CE-RD-01)}
\label{tab:exp-ce-rd-01}
\begin{tabular}{@{}llrrr@{}}
\toprule
Гейт / Наблюдаемое & Теоретическая формула & Значение (Lean 4 / Python) & Допуск & Статус \\
\midrule
Тепловой барьер $D_{\mathrm{thermal}}(13)$ & $1172\cdot 13^2 + 8192$ & $206\,260$ & $0$ & PASS \\
Тепловой барьер $D_{\mathrm{thermal}}(64)$ & $1172\cdot 64^2 + 8192$ & $4\,808\,704$ & $0$ & PASS \\
Когерентный пол $D_{\mathrm{coherent}}(64)$ & $|A|=139 \implies \varepsilon$ & $8\,192$ & $0$ & PASS \\
Выигрыш барьера $(K_{13}, K_{64}, K_{\max})$ & $\lfloor D / 8192\rfloor$ & $(25\times,\ 587\times,\ 833\times)$ & $0$ & PASS \\
Абляция защиты $D>0$ на $t=1\to 2$ & $N_{xy}(2)\colon 26 \to 24$ & $D_{\mathrm{raw}}=-113\,826,\ 26\ne 24$ & $0$ & PASS \\
\bottomrule
\end{tabular}
\end{table}

\subsubsection*{Анализ, расчёт погрешности и выводы}
Все вычисления выполнены в точной целочисленной арифметике $\mathbb{Z}$ без погрешности округления ($\Delta = 0$, тип~B) и проверены ядром Lean~4 без использования аксиом (`does not depend on any axioms`).

\textbf{Вывод.}
Перевод $12$ лучей кубооктаэдра во второй квадрант фазового кольца ($|A_{xy}|\ge 139$ тиков) компенсирует собственный барьер узла $+128\mu^2$ отрицательной продольной голономией $\zeta_{\mathrm{real}}\le -132\mu^2$ и снижает порог топологического пробоя смены ветви в $25\text{--}587$ раз без теплового разогрева числителя.
\textbf{Статус записи:} опыт закрыт по табл.~\ref{tab:exp-ce-rd-01}.

\end{labrecord}

\begin{labrecord}{CE-RD-02 · Квантовый фазовый храповик \texorpdfstring{$-1$}{-1} тик на доменной стенке (\texttt{Int.fdiv}) и топологический сброс \texorpdfstring{$512\,n_f$}{512 n\_f}}{lab:ce-rd-02}
\label{exp:cerd02}

\subsubsection*{Постановка}
\textbf{Физический вопрос.}
Порождает ли целочисленное деление с остатком вниз $\lfloor u/D\rfloor$ (\texttt{Int.fdiv}) в законе ворот~$\Phi$ неустранимую киральную асимметрию $-1$ тик между противоположными берегами плоской доменной стенки при симметричном по модулю возбуждении $\pm u$, и сохраняется ли строго сумма топологических зарядов $\Delta n_{f_+} + \Delta n_{f_-} = 0$ при каждом сбросе $512$ тиков связности через общее ребро?

\textbf{Что проверяем.}
\begin{enumerate}[label=\textbf{(\alph*)},leftmargin=*,itemsep=0.25em]
\item Асимметрию импульсов $\Phi_+ = +76$ и $\Phi_- = -77$ тиков ($\Phi_+ + \Phi_- = -1$ тик) на плоской доменной стенке $\mathrm{IV}(3)=263$ при $t=0$ (\texttt{iv263\_\allowbreak flat\_\allowbreak domain\_\allowbreak wall\_\allowbreak asymmetry\_\allowbreak t0}).
\item Полное исчезновение кирального эффекта ($\Phi_+^{\mathrm{tdiv}} + \Phi_-^{\mathrm{tdiv}} = +76 + (-76) = 0$) при замене деления вниз \texttt{Int.fdiv} на усечение к нулю \texttt{Int.tdiv} (\texttt{ablation\_\allowbreak tdiv\_\allowbreak destroys\_\allowbreak domain\_\allowbreak wall\_\allowbreak chiral\_\allowbreak asymmetry}).
\item Точное равенство $(d_1 A)_f = 512\,n_f$ и сохранение заряда $\Delta n_{f_+} + \Delta n_{f_-} = 0$ при перебросе общего ребра на $\Delta k_e = \pm 1$ (\texttt{plaquette\_\allowbreak flux\_\allowbreak eq\_\allowbreak 512\_\allowbreak mul\_\allowbreak winding}, \texttt{shared\_\allowbreak edge\_\allowbreak charge\_\allowbreak conservation}).
\end{enumerate}

\textbf{Рабочая гипотеза.}
Деление с остатком вниз $\lfloor u/D\rfloor$ действует как точный квантовый храповик с шагом $-1$ тик за такт на любой доменной стенке с $u\bmod D\ne 0$, а каждый переход ребра через разрез Паули $\pm 256$ высвобождает или запасает ровно $512$ тиков фазовой связности при нулевом изменении суммарного топологического заряда.

\textbf{Наблюдаемые.}
Импульсы $\Phi_+, \Phi_-$, сумма храповика $\Delta\Phi_{\mathrm{ratchet}} = \Phi_+ + \Phi_-$, результат абляции $\Delta\Phi_{\mathrm{tdiv}}$ и квант сброса потока $\Delta B_f$ (табл.~\ref{tab:exp-ce-rd-02}).

\textbf{Критерий согласия с теорией.}
Точные целочисленные равенства $\Delta\Phi_{\mathrm{ratchet}} = -1$, $\Delta\Phi_{\mathrm{tdiv}} = 0$, $\Delta B_f = 512\,\Delta n_f$ и $\Delta n_{f_+} + \Delta n_{f_-} = 0$ (тип~B, нулевой допуск).

\subsubsection*{Теоретическая часть}
Для любого целого числителя $u = qD + r$ с остатком $0 < r < D$ имеем $\lfloor +u/D\rfloor = q$, тогда как $-u = -(q+1)D + (D-r)$ даёт $\lfloor -u/D\rfloor = -(q+1)$, откуда $\lfloor +u/D\rfloor + \lfloor -u/D\rfloor = -1$ независимо от величины~$q$.

\textbf{Проверка размерностей.}
Фазовые импульсы $\Phi_\pm$, реберная связность $A_e \in (-256,256]$ и плакетный поток $B_f = 512\,n_f$ измеряются в безразмерных тиках кольца $\mathbb{Z}_{512}$ ($[\Phi] = [B_f] = 1\ \text{тик}$), которым на макроуровне~$\Tlayer$ отвечает квант действия $s_0 = h / 512$.

\subsubsection*{Ход эксперимента}
{\sloppy Воспроизведение: \texttt{lake build} в \texttt{lean/} (\texttt{WorldTrajectory.lean}, \texttt{HolonomyMaxwell.lean}) и \texttt{python verify\_principles.py --suite m\_axioms}.\par}

\begin{enumerate}[label=\textbf{Шаг \arabic*:},leftmargin=*,itemsep=0.35em]
\item Вычислить голономии $(\zeta_{\mathrm{real}},\zeta_{\mathrm{imag}}) = (134\,862,\ \pm 70\,980)$ на противоположных сторонах плоской стенки октантов $\mathrm{IV}(3)=263$.
\item Сравнить выход ворот при делении вниз \texttt{Int.fdiv} ($+76$ и $-77$) и при усечении к нулю \texttt{Int.tdiv} ($+76$ и $-76$).
\item Проверить рождение первого вихревого плакета $B_f(2) = 32 + 85 + 196 + 199 = 512$ ($n_f(2) = +1$) в центральном углу~$K_3$ на такте $t=2$.
\end{enumerate}

\begin{table}[htbp]
\centering
\footnotesize
\setlength{\tabcolsep}{4pt}
\caption{Верификация квантового фазового храповика и топологического сброса $512\,n_f$ (CE-RD-02)}
\label{tab:exp-ce-rd-02}
\begin{tabular}{@{}llrrr@{}}
\toprule
Гейт / Наблюдаемое & Теоретическая формула & Значение (Lean 4 / Python) & Допуск & Статус \\
\midrule
Импульсы стенки $(\Phi_+, \Phi_-)$ & $\lfloor\pm 12\,524\,928 / 164\,686\rfloor$ & $(+76,\ -77)$ тиков & $0$ & PASS \\
Шаг фазового храповика $\Phi_+ + \Phi_-$ & $\lfloor +u/D\rfloor + \lfloor -u/D\rfloor$ & $-1$ тик & $0$ & PASS \\
Абляция \texttt{Int.tdiv} ($\Phi_+^{\mathrm{tdiv}} + \Phi_-^{\mathrm{tdiv}}$) & $76 + (-76)$ & $0$ (эффект исчезает) & $0$ & PASS \\
Квант плакетного потока $B_f(t=2)$ & $(d_1 A)_f = 512\,n_f$ & $32+85+196+199 = 512$ & $0$ & PASS \\
Закон сохранения на общем ребре & $\Delta n_{f_+} + \Delta n_{f_-}$ & $0$ (тождественно) & $0$ & PASS \\
\bottomrule
\end{tabular}
\end{table}

\subsubsection*{Анализ, расчёт погрешности и выводы}
Арифметика целочисленного деления с остатком точна ($\Delta = 0$, тип~B); обе теоремы в \texttt{WorldTrajectory.lean} не зависят ни от одной аксиомы Lean~4.

\textbf{Вывод.}
Сдвиг округления вниз $\lfloor u/D\rfloor$ на доменной стенке формирует точный однонаправленный фазовый храповик в $-1$ тик за такт, а переброс ребра через разрез Паули $\pm 256$ квантует запас и сброс вихревого потока порциями ровно по $512$ тиков.
\textbf{Статус записи:} опыт закрыт по табл.~\ref{tab:exp-ce-rd-02}.

\end{labrecord}

\begin{labrecord}{CE-RD-03 · Четырёхканальная голономная томография \texorpdfstring{$\mathcal{M}(x,t,N)$}{M(x,t,N)} и устойчивость геодинамо при \texorpdfstring{$D(x)\ge 8\,192$}{D(x)>=8192}}{lab:ce-rd-03}
\label{exp:cerd03}

\subsubsection*{Постановка}
\textbf{Физический вопрос.}
Выполняются ли для замкнутого $4$-канального оператора измерения макроприбора $\mathcal{M}(x,t,N) = \mathcal{B}_N[\mathcal{O}(\Psi^t)](x)$ при любом конечном окне $N\ge 1$ точные тождества голономной томографии ($\delta K_{xy} = \mu_x\mu_y|A_{xy}|$, $J_{xy}+J_{yx}=0$, $(d_1\mathcal{B}_N[A])_f = 512\,\mathcal{B}_N[n_f]$, $\operatorname{div}_{\mathrm{cube}}\mathcal{B}_N[n_f]\equiv 0$) и строгая положительность оптического знаменателя $D(x,t) \ge 8\,192 > 0$, исключающая сингулярность $g_{00}^{(N)}\to 0$ в сжатых ядрах планет?

\textbf{Что проверяем.}
{\sloppy
\begin{enumerate}[label=\textbf{(\alph*)},leftmargin=*,itemsep=0.25em]
\item Замкнутую линейную формулу для чётного дефицита Бома $\delta K_{xy} = 128\mu_x\mu_y - \mu_x\mu_y C(d\varphi) = \mu_x\mu_y|\mathrm{wrap}_{512}(d\varphi)|$ и антисимметрию тока $J_{xy}+J_{yx}=0$ на всех $512$ углах кольца (\texttt{even\_\allowbreak holonomy\_\allowbreak bohm\_\allowbreak deficit\_\allowbreak eq\_\allowbreak abs\_\allowbreak connection}, \texttt{odd\_\allowbreak holonomy\_\allowbreak current\_\allowbreak pairwise\_\allowbreak conservation\_\allowbreak all\_\allowbreak 512}).
\item Точную коммутацию оператора прибора~$\mathcal{B}_N$ с уравнениями Максвелла--Гаусса--Фарадея при любом конечном~$N\ge 1$ (\texttt{instrument\_\allowbreak convolution\_\allowbreak plaquette\_\allowbreak flux\_\allowbreak exact}, \texttt{instrument\_\allowbreak convolution\_\allowbreak gauss\_\allowbreak monopole\_\allowbreak free\_\allowbreak exact}).
\item Строгую нижнюю границу $D(x,t) > 0$ ($D\ge 8\,192$) и нулевую дивергенцию куба $\max|\operatorname{div}_{\mathrm{cube}} n_f| = 0$ на $86$ вихревых плакетах 3D-решётки $\mathrm{IV}(3)=263$ (\texttt{gravity\_\allowbreak denominator\_\allowbreak strictly\_\allowbreak positive\_\allowbreak floor\_\allowbreak 8192}, \texttt{closed\_\allowbreak form\_\allowbreak unified\_\allowbreak 4channel\_\allowbreak measurement\_\allowbreak iv263\_\allowbreak t2}).
\end{enumerate}\par}

\textbf{Рабочая гипотеза.}
На любом конечном масштабе прибора $N\ge 1$ все четыре канала измерения $\mathcal{M}(x,t,N)$ связаны безынтегральными алгебраическими тождествами с нулевой погрешностью ($\Delta \equiv 0$).

\textbf{Наблюдаемые.}
Погрешность линейной реконструкции дефицита Бома $\Delta_{\mathrm{Bohm}}$, невязка тока $\Delta_J = \max|J_{xy}+J_{yx}|$, максимальная дивергенция куба $\max|\operatorname{div}_{\mathrm{cube}}\mathcal{B}_N[n_f]|$, полный вихревой поток $\sum |B_f(t=2)|$ и минимум знаменателя $\min D(x,t)$ (табл.~\ref{tab:exp-ce-rd-03}).

\textbf{Критерий согласия с теорией.}
Точные равенства $\Delta_{\mathrm{Bohm}} = 0$, $\Delta_J = 0$, $\max|\operatorname{div}_{\mathrm{cube}}\mathcal{B}_N[n_f]| = 0$, $\sum |B_f(2)| = 44\,032$ тика и $\min D \ge 8\,192 > 0$ (тип~B).

\subsubsection*{Теоретическая часть}
Поскольку $C(d\varphi) = 128 - |\mathrm{wrap}_{512}(d\varphi)|$, сумма $C(d\varphi) + |\mathrm{wrap}_{512}(d\varphi)| \equiv 128$ позволяет однозначно разделять плотностной множитель $\mu_x\mu_y$ и фазовый разрыв $|A_{xy}|$ по совместному замеру продольной задержки~$D$ и дисперсионного дефицита~$\delta K_{xy}$.

\textbf{Проверка размерностей.}
Оператор окна прибора $W_N(z) = 12^{-2N}\#\{\gamma\in O_2^{2N}\}$ безразмерен и нормирован на единицу ($\sum_z W_N(z) = 1$), поэтому все сглаженные наблюдаемые $\mathcal{B}_N[F]$ сохраняют точную физическую размерность исходных решёточных полей~$F$.

\subsubsection*{Ход эксперимента}
{\sloppy Воспроизведение: \texttt{lake build} в \texttt{lean/} (\texttt{UnifiedOperators.lean}, \texttt{HolonomyMaxwell.lean}, \texttt{TrajectoryIV263.lean}) и \texttt{python verify\_principles.py --suite m\_axioms}.\par}

\begin{enumerate}[label=\textbf{Шаг \arabic*:},leftmargin=*,itemsep=0.35em]
\item Проверить тождества $128 - C(r) = |r|$ и $S(r) + S((512-r)\bmod 512) = 0$ по всем $512$ остаткам кольца $\mathbb{Z}_{512}$.
\item Вычислить 3D-эволюцию $\symCapitalPsi^t = g^t(\mathrm{IV}(3)=263)$ при $t=0,1,2$ и измерить полный вихревой поток $512\times(26_{xy}+26_{yz}+34_{xz})$ и дивергенцию во всех элементарных кубах.
\item Проверить выполнение коммутационных тождеств прибора~$\mathcal{B}_N$ для произвольного конечного списка весов~$w_i\in\mathbb{Z}$.
\end{enumerate}

\begin{table}[htbp]
\centering
\footnotesize
\setlength{\tabcolsep}{4pt}
\caption{Верификация 4-канальной голономной томографии и регулярности ядра (CE-RD-03)}
\label{tab:exp-ce-rd-03}
\begin{tabular}{@{}llrrr@{}}
\toprule
Гейт / Наблюдаемое & Теоретическая формула & Значение (Lean 4 / Python) & Допуск & Статус \\
\midrule
Линейная формула Бома $\delta K_{xy}$ & $\mu_x\mu_y(128 - C) = \mu_x\mu_y|A|$ & $\Delta_{\mathrm{Bohm}} = 0$ ($\forall\mu,d\varphi\in\mathbb{Z}$) & $0$ & PASS \\
Антисимметрия тока $J_{xy}+J_{yx}$ & $\mu_x\mu_y(S(r)+S(-r))$ & $0$ (на всех $512$ углах) & $0$ & PASS \\
Поток вихревого ядра $\sum |B_f(2)|$ & $512\times(26+26+34)$ & $44\,032$ тика ($86$ плакетов) & $0$ & PASS \\
Дивергенция под прибором $\mathcal{B}_N$ & $\operatorname{div}_{\mathrm{cube}}\mathcal{B}_N[n_f]$ & $0$ ($\forall N\ge 1$) & $0$ & PASS \\
Нижняя грань импеданса $\min D(x,t)$ & $\max(D_{\mathrm{raw}}, \varepsilon)$ & $\ge 8\,192 > 0$ ($g_{00}^{(N)} > 0$) & $0$ & PASS \\
\bottomrule
\end{tabular}
\end{table}

\subsubsection*{Анализ, расчёт погрешности и выводы}
Все пять операторных тождеств доказаны в Lean~4 для произвольных целочисленных входов без аппроксимаций ($\Delta = 0$, тип~B).

\textbf{Вывод.}
Замкнутая 4-канальная формула $\mathcal{M}(x,t,N)$ обеспечивает точную алгебраическую реконструкцию внутренних профилей плотности и фазовой связности при любом конечном окне прибора $N\ge 1$ и гарантирует отсутствие сингулярностей в сверхплотных ядрах.
\textbf{Статус записи:} опыт закрыт по табл.~\ref{tab:exp-ce-rd-03}.

\end{labrecord}
\end{labseries}

